\section{Related Work}
\label{sec:related}

\subsection{Code Translation}
Code translation facilitates cross-platform interoperability and promotes code reuse across diverse programming environments, playing a crucial role in legacy system modernization~\cite{lu2021codexglue,puri2021codenet}.

Early approaches evolved from rule-based systems with handcrafted transformation rules to Statistical Machine Translation (SMT) for modeling common code patterns~\cite{nguyen2013lexical,chen2018tree}, and subsequently to neural network models, though early RNNs struggled with complex syntactic structures~\cite{dong2016language,yin2017syntactic,yin2018tranx}.

LLMs have revolutionized code translation with unprecedented capabilities, though challenges remain in translation accuracy and repository-specific context handling~\cite{zhu2022xlcost,rithy2022xtest}.
Recent work has addressed these challenges through various approaches: benchmarking LLM limitations in complex scenarios~\cite{jiao2023evaluation}, context-enriched prompts~\cite{pan2024lost}, test case feedback for iterative refinement~\cite{yang2024exploring}, consistent prompt design~\cite{macedo2024exploring}, and comprehensive repository-aware evaluation~\cite{ou2024repository}. Despite these advances, repository-aware contextual understanding remains largely unexplored. Our work addresses this limitation through a multi-agent framework specifically designed for repository-aware code translation.

\subsection{LLM-based Multi-Agent Systems}
LLM-based multi-agent systems have emerged as a promising research direction, advancing several areas. 
Infrastructure development focuses on frameworks for efficient agent coordination and task management~\cite{gong2023mindagent,chen2023agentverse,zhang2023building,hong2023metagpt,zhang2024proagent}. 
Benchmark development evaluates multi-agent performance in dynamic environments~\cite{chen2024llmarena,dong2024villageragent}. 
Large-scale social simulations explore modeling complex societal behaviors~\cite{al2024project,li2022evolution,de2004epidemiology}. 
Domain-specific applications demonstrate practical effectiveness in specialized scenarios, including software engineering tasks~\cite{d2024marg,zhang2024autocoderover,yang2024swe}.
Our work contributes to this research by introducing a specialized multi-agent framework for repository-aware code translation. Unlike general-purpose systems, \toolname leverages domain-specific knowledge and specialized agents to address the unique challenges of cross-language code translation in real-world software repositories.